%% notes-tex/multimodal/043-training-design/sections/3-model.tex
%% SOURCE: torchcell/models/cell_model_043.py (CompositeToken, PositionTerm, CellTrunk043,
%% MarkOperator043, CellModel043), torchcell/models/readout_heads_043.py,
%% torchcell/models/genotype_tokens_043.py (ChangeType, CHANGE_PRECEDENCE),
%% torchcell/models/equivariant_cell_graph_transformer.py (GraphRegularizedTransformerLayer),
%% torchcell/data/representation_043.py (GapKind, PositionGap), train_043.py (EXP_TYPE_PATHS,
%% build_model), conf/train_043_default.yaml; parameter count and table rows from the
%% job 6136 log.

\section{Model, step by step\secstatus{todo}}
\label{sec:model}

The model keeps the manuscript's factorization $\hat f_\theta = \mathcal{R}_\phi \circ
\mathcal{T}_\psi \circ F_\theta$ (Methods, ``Problem setup'') and changes what enters it:
the input token of a gene is a function of the reference context, the perturbation tokens
are typed marks, and the decoder is two shared heads with the dataset entering only at
their input and output. Shapes are given for one record unless a batch axis is written;
$C$ is the number of distinct contexts in a batch.

\subsection{The gene token\secstatus{todo}}
\label{sec:token}

\notesec{Step 1, the four frozen blocks}
For context $c$ and gene $i$, each block $q \in \mathcal{Q}$ holds either a key or a typed
gap $\kappa \in \mathcal{K}$. A key indexes a frozen fp16 vector $x^{(q)}_{c,i} \in
\mathbb{R}^{d_q}$: the FUDT upstream encoder on the 1{,}003 bp window ($d_{\mathrm{up}} =
768$), the FUDT downstream encoder on the 300 bp window ($d_{\mathrm{down}} = 768$),
ProtT5-XL on the protein ($d_{\mathrm{prot}} = 1024$) and CaLM on the coding sequence
($d_{\mathrm{calm}} = 768$). Each block is normalized and projected separately,
\begin{equation}
\tilde x^{(q)}_{c,i} =
\begin{cases}
W^{(q)} \operatorname{LN}^{(q)}\!\big(x^{(q)}_{c,i}\big) + b^{(q)} & \text{if block } q \text{ of gene } i \text{ has a key in } c,\\[2pt]
\mathbf{E}^{\mathrm{gap}}_{q,\kappa} & \text{if it carries the gap } \kappa,
\end{cases}
\qquad W^{(q)} \in \mathbb{R}^{d \times d_q},
\label{eq:block}
\end{equation}
so a gap is a learned vector of its own kind and block and is never filled with another
sequence's vector. The vectors stay frozen; $\operatorname{LN}^{(q)}$, $W^{(q)}$, $b^{(q)}$
and $\mathbf{E}^{\mathrm{gap}} \in \mathbb{R}^{(4 \cdot 12) \times d}$ train. The overnight
runs held 929{,}348, 326{,}255, 329{,}430 and 514{,}189 rows of the four tables on the card,
3.39 GB (job 6136 log).

\notesec{Step 2, the composite}
The four projections are summed,
\begin{equation}
\tilde x_{c,i} = \sum_{q \in \mathcal{Q}} \tilde x^{(q)}_{c,i} \in \mathbb{R}^{d}.
\label{eq:composite}
\end{equation}

\notesec{Step 3, the position term}
Gene $i$ lies on chromosome $\operatorname{chr}(i) \in \{0, \ldots, 16\}$ (0 is the
mitochondrion) at relative midpoint $\xi_i \in [0,1]$ of the reference. With $\nu^{\mathrm{chr}}_{c,i}$
the copy number of that chromosome relative to the basal ploidy (1 except on aneuploid
isolate chromosomes) and $\nu^{\mathrm{gene}}_{c,i}$ the gene's copy number,
\begin{equation}
\operatorname{pos}_{c,i} =
\nu^{\mathrm{chr}}_{c,i}\, \mathbf{E}^{\mathrm{chr}}_{\operatorname{chr}(i)}
+ W_{\xi} \operatorname{Four}(\xi_i) + b_{\xi}
+ \mathbf{w}_{\nu}\big(\log_2(\nu^{\mathrm{gene}}_{c,i} + 1) - 1\big) + \mathbf{b}_{\nu}
+ \mathbf{E}^{\mathrm{pgap}}_{\kappa'_{c,i}},
\label{eq:pos}
\end{equation}
with $\mathbf{E}^{\mathrm{chr}} \in \mathbb{R}^{19 \times d}$ (17 chromosomes and two rows for
off-graph tokens, an accessory ORF and the reporter construct), $W_{\xi} \in \mathbb{R}^{d
\times 16}$ and $\mathbf{w}_\nu, \mathbf{b}_\nu \in \mathbb{R}^{d}$. The copy feature is 0 at
one copy and $-1$ at zero copies. An unknown copy number (a ploidy the source did not report,
a value the pangenome matrix does not release) drops its term and selects the learned
position-gap row $\kappa'_{c,i}$ of $\mathbf{E}^{\mathrm{pgap}} \in \mathbb{R}^{4 \times d}$
(row 0 when nothing is unknown).

\notesec{Step 4, the input token}
The manuscript's node feature $x_i = \operatorname{Embed}(g_i)$ becomes context dependent,
\begin{equation}
x_{c,i} = \tilde x_{c,i} + \operatorname{pos}_{c,i},
\qquad
H^{(0)}_c = (x_{\mathrm{CLS}}, x_{c,1}, \ldots, x_{c,N}) \in \mathbb{R}^{(N+1) \times d}.
\label{eq:token}
\end{equation}
For S288C every slot has its reference keys; under route A an isolate's slots carry the
isolate's own alleles, with the genes outside the released allele set and the genes the
record calls absent as gaps.

\subsection{The trunk\secstatus{todo}}
\label{sec:trunk}

\notesec{Step 5, the layers}
$F_\theta$ is four layers of width $d = 144$ with $h = 9$ heads of width $d_k = 16$. The
layer is the manuscript's Eqs.\ for the queries, keys, values and attention map, but
post-normalized, where the manuscript writes the pre-normalized form,
\begin{equation}
\hat H^{(\ell)} = \operatorname{LN}\!\Big(H^{(\ell-1)} + \operatorname{Drop}\big(\big[O^{(\ell,1)} \Vert \cdots \Vert O^{(\ell,h)}\big] W_O^{(\ell)}\big)\Big),
\qquad
H^{(\ell)} = \operatorname{LN}\!\Big(\hat H^{(\ell)} + \operatorname{Drop}\big(\operatorname{FFN}^{(\ell)}(\hat H^{(\ell)})\big)\Big).
\label{eq:layer}
\end{equation}

\notesec{Step 6, the graph mask}
In layer $\ell = 2$ (\texttt{mask\_layers: [1]}, counted from zero) head $a$ may attend only
along graph $k(a)$; the other three layers are unmasked. With $A^{(k)}$ symmetrized,
\begin{equation}
\alpha^{(2,a)} = \operatorname{softmax}\!\Big(\frac{Q^{(2,a)} (K^{(2,a)})^{\top}}{\sqrt{d_k}} + \Gamma^{(a)}\Big),
\qquad
\Gamma^{(a)}_{ij} =
\begin{cases}
0 & \text{if } A^{(k(a))}_{ij} = 1,\ A^{(k(a))}_{ji} = 1,\ i = j, \text{ or } i \text{ or } j \text{ is the CLS},\\
-\infty & \text{otherwise},
\end{cases}
\label{eq:mask}
\end{equation}
with $k(0), \ldots, k(8)$ the physical, regulatory and TFLink graphs and the STRING 12.0
neighborhood, fusion, co-occurrence, coexpression, experimental and database channels.
The objective carries no $\Omega_k$ term: 043 takes the manuscript's $\lambda \to \infty$
limit of graph-regularized attention at one layer and leaves the rest free.

\notesec{Step 7, one encode per context}
Each distinct context of a batch is encoded once, $H_c = F_\theta(G_c) \in \mathbb{R}^{(N+1)
\times d}$, a $[C, N+1, d]$ tensor, and record $b$ reads $H_{c(b)}$. The trunk trains, so
nothing is cached across batches.

\subsection{Marks and the perturbation operator\secstatus{todo}}
\label{sec:operator}

\notesec{Step 8, change types}
Each gene of a genotype has one change type in $\mathcal{C}$. When the source gives a gene
several, the first in the order absent, deletion, loss of function, variant, present,
promoter wins: a gene that is gone outranks a call on its sequence, which outranks a
variant that is still expressed. A mark is a token the operator attends to: a deleted gene,
a gene a segregant inherits from a non-BY parent (a variant, with that parent as $o_t$), the
HO deletion of an isolate under route A, an isolate's present accessory ORF (off-graph), and
Vaishnav's reporter, an off-graph token whose upstream key is the variant's 1{,}003 bp
construct window. Under route A the isolate's variants live in its context, not in its
marks. The overnight runs kept at most 1{,}024 marks per record (\texttt{data.mark\_cap}:
absent, deletion and loss-of-function marks first, then a fixed per-record random subset),
a smoke-time truncation the design does not have; N'Guessan 2025 records carry a median of
2{,}910 marks (job 6120).

\notesec{Step 9, the mark key}
Mark $t$ of record $b$ has the composite and position of its gene in the record's context
(or of its off-graph token, with the replicon row of $\mathbf{E}^{\mathrm{chr}}$ and the
Fourier term zeroed), plus its change type and allele parent,
\begin{equation}
\tilde{\mathbf{p}}_t = \tilde x_{t} + \operatorname{pos}_{t} + \mathbf{t}_{\tau_t} + \mathbf{E}^{\mathrm{par}}_{o_t},
\qquad
\mathbf{p}_t = W_p \operatorname{LN}(\tilde{\mathbf{p}}_t) + b_p + \mathbf{1}[\text{mark } t \text{ is on the graph}]\; h_{c(b),\,j_t},
\label{eq:markkey}
\end{equation}
where $j_t$ is the marked gene and $\mathbf{E}^{\mathrm{par}}_0 = 0$ is ``no parent''. This is
the manuscript's token $\mathbf{p}_t = r(e_t) + \mathbf{t}_{\tau_t} + m_t \mathbf{w}$ with two
changes: the projected mark token is added for in-cell genes as well, and there is no
magnitude term. $\mathbf{P}_b \in \mathbb{R}^{M_b \times d}$, padded to $[B, \max_b M_b, d]$.

\notesec{Step 10, cross-attention with a null sink}
Every token of the context, the CLS included, is a query. Each head adds a learned null key
$\mathbf{k}_{\varnothing}$ with value zero, so a record without marks attends only to it and a
record with one mark does not see a degenerate softmax,
\begin{equation}
\beta^{a}_{i,t} = \operatorname*{softmax}_{t \in \{\varnothing, 1, \ldots, M_b\}} \frac{(h_i W^{a}_Q) \cdot \mathbf{k}^{a}_t}{\sqrt{d_k}},
\quad
\mathbf{k}^{a}_t = \mathbf{p}_t W^{a}_K,
\quad
\Delta h_i = \Big[\textstyle\sum_{t \ge 1} \beta^{1}_{i,t}\, \mathbf{p}_t W^{1}_V \,\Big\Vert \cdots \Big\Vert\, \sum_{t \ge 1} \beta^{h}_{i,t}\, \mathbf{p}_t W^{h}_V\Big] W_O,
\label{eq:opattn}
\end{equation}
for $i \in \{\mathrm{CLS}, 1, \ldots, N\}$; the null term contributes weight but no value.

\notesec{Step 11, the ReZero residual}
Where the manuscript's operator normalizes, 043 gates,
\begin{equation}
h'_i = h_i + \zeta_{\mathrm{attn}} \operatorname{Drop}(\Delta h_i),
\qquad
h^{\mathrm{pert}}_i = h'_i + \zeta_{\mathrm{ffn}} \operatorname{Drop}\big(\operatorname{FFN}(h'_i)\big),
\qquad
H_{\mathrm{pert}} = \mathcal{T}_\psi(H_{c(b)}, p_b) \in \mathbb{R}^{(N+1) \times d},
\label{eq:rezero}
\end{equation}
with $\zeta_{\mathrm{attn}} = \zeta_{\mathrm{ffn}} = 0$ at initialization, so at step 0 the
operator is the identity and $H_{\mathrm{pert}} = H_{c(b)}$. The record-level summary of the
marks is
\begin{equation}
\bar{\mathbf{k}}_b = \frac{1}{\max(M_b, 1)} \sum_{t=1}^{M_b} \mathbf{p}_t \in \mathbb{R}^{d}.
\label{eq:summary}
\end{equation}

\subsection{Conditioning tokens and heads\secstatus{todo}}
\label{sec:heads}

\notesec{Step 12, the experiment-type token}
Each dataset reads one phenotype class of the schema, and the class enters as the sum of
one learned row per node of its path,
\begin{equation}
\mathbf{c}^{\mathrm{type}}_b = \sum_{\chi \in \operatorname{path}(\operatorname{class}(s(b)))} \mathbf{E}^{\mathrm{type}}_{\chi} \in \mathbb{R}^{16},
\label{eq:typetoken}
\end{equation}
where a path runs root, group, leaf: phenotype, expression, then the microarray (Kemmeren),
RNA-seq (Caudal) or single-cell (Boocock, N'Guessan) class; phenotype, protein, then protein
abundance (Messner, Zelezniak, Muenzner) or turnover (Muenzner turnover); phenotype,
promoter, pooled sort-seq (Vaishnav). The path is the minimal differential subset of the
ontology DAG for these six leaves: it keeps the one level that separates them into the three
readout groups the heads share, so two leaves under one group share that group's row and
differ by their own. The group nodes are modeling-side groups, not schema classes
(\file{train_043.py}, \texttt{EXP\_TYPE\_PATHS}); the root is on every path and acts as a
shared offset. The flat alternative, one row per leaf, is \texttt{exp\_type\_mode: flat}
(Sec.~\ref{sec:open}).

\notesec{Step 13, the conditioning token}
The dataset and the environment are flat lookups,
\begin{equation}
\mathbf{c}_b = \big[\mathbf{c}^{\mathrm{type}}_b ;\ \mathbf{E}^{\mathrm{data}}_{s(b)} ;\ \mathbf{E}^{\mathrm{env}}_{\varepsilon(b)}\big] \in \mathbb{R}^{16 + 8 + 16 = 40},
\label{eq:cond}
\end{equation}
over the nine datasets and nine environments of the prepared store.

\notesec{Step 14, the gene head}
Every gene-level dataset shares one head that reads the gene's own perturbed state, the
manuscript's equivariant readout with the conditioning token appended,
\begin{equation}
\hat{\mathbf{o}}_{b,i} = \operatorname{MLP}_{\mathrm{gene}}\big([h^{\mathrm{pert}}_{b,i} ;\ \mathbf{c}_b]\big) + \mathbf{b}^{\mathrm{gene}}_{s(b)} \in \mathbb{R}^{n_{\mathrm{out}}},
\qquad i = 1, \ldots, N,
\label{eq:genehead}
\end{equation}
with $\operatorname{MLP}_{\mathrm{gene}}: \mathbb{R}^{184} \to \mathbb{R}^{128} \to
\mathbb{R}^{n_{\mathrm{out}}}$, a $[B_{\mathrm{gene}}, N, n_{\mathrm{out}}]$ output.

\notesec{Step 15, the record head}
Vaishnav's one value per promoter is read from the perturbed CLS, the mark summary and the
conditioning token,
\begin{equation}
\hat{\mathbf{o}}_{b} = \operatorname{MLP}_{\mathrm{rec}}\big([h^{\mathrm{pert}}_{b,\mathrm{CLS}} ;\ \bar{\mathbf{k}}_b ;\ \mathbf{c}_b]\big) + \mathbf{b}^{\mathrm{rec}}_{s(b)} \in \mathbb{R}^{n_{\mathrm{out}}},
\label{eq:rechead}
\end{equation}
with $\operatorname{MLP}_{\mathrm{rec}}: \mathbb{R}^{328} \to \mathbb{R}^{128} \to
\mathbb{R}^{n_{\mathrm{out}}}$; it replaces the manuscript's perturbed-gene mean in the
interaction head by the mark-key mean of Eq.~\ref{eq:summary}.

\notesec{The per-dataset term is a bias}
The design calls for a per-dataset affine at the head, $a_s f(\cdot) + \mathbf{b}_s$.
\texttt{CellModel043} builds the gene head with \texttt{per\_dataset\_affine=False} and adds
$\mathbf{b}_s$ alone, so the scale $a_s$ is fixed at 1 in every run so far. Under HL-Gauss
$\mathbf{b}_s$ adds to the 32 logits, a learned prior histogram per dataset. The dataset
id therefore enters in two places, both at the heads: $\mathbf{E}^{\mathrm{data}}$ inside
$\mathbf{c}_b$ and the bias. It never reaches the trunk or the operator.
